\section{Introduction}
\label{sec:introduction}

A response from a server to a client is, more often than not, a list of
records: the rows of a report, the products of a catalogue page, the events of
a time window. JSON~\cite{rfc8259} writes such a list with every field name in
every record and every number as decimal text. The binary formats that replace
it in practice, Protocol Buffers~\cite{protobuf}, MessagePack~\cite{msgpack}
and CBOR~\cite{rfc8949}, remove the names or shorten them, and write numbers
in binary, but they keep the row-by-row layout: a thousand records carry a
thousand copies of each field's tag, and values of one field, which resemble
each other, are separated by values of the others, which do not.

Columnar formats solved this for analytical storage long ago. Storing the
values of one column together removes the per-row tags, and lets each column
be encoded with lightweight transforms such as delta, frame of reference and
bit-packing~\cite{zukowski2006pfor,lemire2015simdbp128}, as
Parquet~\cite{parquet} and, more recently, BtrBlocks~\cite{kuschewski2023btrblocks}
and FastLanes~\cite{afroozeh2023fastlanes} do for files, and as
Arrow~\cite{arrow} does for in-memory interchange without compression. These
formats are designed for tables of millions of rows. Their framing, schema and
metadata are a fixed cost that a message of a few hundred bytes cannot
absorb, they do not carry a document that is not a table, and they come with
a schema language or a builder API that the application's own types have to
be mapped to.

This paper presents \colbin{}, a \emph{hybrid} serialization format: a
message is written in a row format, and the columnar layout is applied
inside it, where it pays, and nowhere else. The unit of decision is the
value, at encode time: a slice of structs whose fields are scalars or
strings is written as a list of rows when it is short, and as a table of
columns from eight rows up. The rest of the message, the fields of the root,
nested structs and maps, uses the row format, with one-byte field headers
and no variable-length integers. A three-field configuration document and a
million-row table are therefore the same format, and the encoder needs no
hint about which it is writing. The two halves are evaluated apart: the
single records of the evaluation measure the row format alone, and the
collections measure the columns.

The second design constraint is setup cost. The schema of a \colbin{} message
is the Go struct itself: field numbers come from optional struct tags or are
derived from field names, there is no interface description file, no compiler
step and no generated type. This is the niche of Go's
\fmt{gob}~\cite{gob,pike2011gobs}, and the reason a format with the size of a
schema-compiled one is worth having without the compilation. A message can
also carry its schema, in which case any reader renders it as JSON with no
type at hand.

\paragraph{What it is used for.} \colbin{} is a general-purpose format: any
place where Protocol Buffers or MessagePack would serve is a candidate. It
was built for one application, a business system developed by the author,
and its design follows the places where that system had JSON. They are of three kinds, and the
evaluation is organised by them.
\begin{itemize}
\item \emph{Bulk tables}: archives and backups of whole tables, written once
  and read rarely, where size is the cost.
\item \emph{Tables of web size}: the responses a server sends to a web
  client, a page of records at a time, and the cache of those responses.
\item \emph{Single records}: a record stored as one binary value in a
  key-value store, in place of a JSON document or one attribute per field,
  a small object stored in a binary column of a relational table, where a
  JSON or text column would otherwise hold it, and a message between two
  services.
\end{itemize}
The row format's narrow field header, one byte for a type of up to sixteen
fields (Section~\ref{sec:format}), exists for the last kind: the small
objects that end up as a field of a larger record are, in a database,
usually a JSON string, and most of such a string is field names and
punctuation. A format that is to replace JSON in all three places has to
be cheap for an object of three fields and compact for a table of a
million rows, which is the choice per value described above.

\paragraph{Contributions.}
\begin{itemize}
\item A hybrid message format in which a slice of structs is a list of rows
  or a columnar table, decided per value, with integer columns encoded by a
  per-column transform and bit-packed in byte-aligned blocks of 128 values,
  and float columns encoded as decimals where they are decimals
  (Section~\ref{sec:format}).
\item A row format for everything that is not a table: one byte of header
  per field for types of up to sixteen fields, zero values not written, and
  every size read with a single load rather than a variable-length loop.
\item A self-describing mode: a message can embed its schema, a key, a name
  and an operation code per field, and any reader then renders it as JSON,
  or decodes it into another type by field name, without the producer's
  type. Its cost is measured against the message it describes
  (Section~\ref{sec:composition}).
\item An evaluation in Go against \result{formats.baselines} formats and
  implementations (JSON, \fmt{gob}, CBOR, MessagePack, Protocol Buffers with two
  implementations, FlatBuffers, Arrow IPC, Avro and Cap'n Proto), raw and
  under gzip and Zstandard, on \result{cases.total} cases from four
  public sources, none of them designed by the author
  (Sections~\ref{sec:methodology} and~\ref{sec:results}).
\item Ablations of the two parameters the layout rests on: the row count at
  which a slice becomes a table, and the choice of column transform
  (Section~\ref{sec:ablations}).
\item An account of where the format loses or ties, and of what the
  evaluation does not cover (Section~\ref{sec:limitations}).
\item A reproducible artifact: datasets pinned by checksum, the exact subset
  rules, and one command that regenerates every table and figure of this
  paper (Section~\ref{sec:availability}).
\end{itemize}

The evaluation compares \colbin{} with each other format directly, on the
cases both carry, and reports the comparison two ways: by dataset, and by
the three kinds of message above. A response to a web client bounds the
size of a message, so a table of web size is one of at most $2^{14}$ rows
and a bulk table one past that bound; a single record is a document that
holds no table.
